%%=============================================================================
%% Reverse engineering van TPPBAMT
%%=============================================================================

\chapter{\IfLanguageName{dutch}{Reverse engineering}{Reverse engineering}}%
\label{ch:reverse-engineering}

In dit hoofdstuk wordt \emph{Component~1} van het migratieframework, de analysefase, toegepast op de
CA~Gen-ActionBlock \texttt{AB\_APP\_GBP\_MAT\_BAF}. Het doel is een volledig analysedossier op te bouwen dat als contract dient
voor de PL/I-implementatie in Hoofdstuk~\ref{ch:poc} en voor de validatie in Hoofdstuk~\ref{ch:validatie}.
Het hoofdstuk volgt de structuur van \autoref{sec:component-analyse}, met als bijkomende doelstelling het
framework zelf te toetsen aan een re\"ele casus.

\section{Encyclopedie en versiebeheer}
\label{sec:re-encyclopedie}

In een eerste stap werd het businessmodel waaronder \texttt{AB\_APP\_GBP\_MAT\_BAF} valt, lokaal beschikbaar gemaakt via
een \emph{Check Out} uit de host-encyclopedie en geopend in de CA~Gen-toolset.

\section{View-inventaris}
\label{sec:re-views}

Per Action Block binnen \texttt{AB\_APP\_GBP\_MAT\_BAF} werd een view-inventaris opgesteld volgens de typologie uit
\autoref{sec:Views} en de template uit \autoref{tab:view-inventaris}. \autoref{tab:tppbamt-views-overzicht}
geeft een sterk vereenvoudigd overzicht van de hoofdcategorie\"en views die in de module voorkomen. 
%De volledige inventaris is opgenomen in Bijlage~\ref{ch:bijlage-views}.

\begin{table}[h]
\centering
\caption[Overzicht view-categorie\"en in \texttt{AB\_APP\_GBP\_MAT\_BAF}]{Overzicht van view-categorie\"en in \texttt{AB\_APP\_GBP\_MAT\_BAF}.}
\label{tab:tppbamt-views-overzicht}
\begin{tabular}{|l|c|p{0.45\linewidth}|}
\hline
\textbf{Categorie} & \textbf{\#} & \textbf{Typische rol} \\
\hline
Import View & 11 & Materiaaldetails. \\
\hline
Export View & 2 & Statusvelden, exitstate, samenvattende kentallen. \\
\hline
Locals & 2 Work Views & Tijdelijke berekeningen, tellers, ... \\
\hline
Entity Actions & 1 & SQL velden van een RO. \\
\hline
Work View & 10+ & Niet-gepersisteerde tussenresultaten en flags. \\
\hline
Group View & 10+ & Werkzones. \\
\hline
\end{tabular}
\end{table}

\section{Control-flow en USE-graaf}
\label{sec:re-control-flow}

De control-flow van \texttt{AB\_APP\_GBP\_MAT\_BAF} werd ontleed. De hoofdlus itereert over
de planbare materialen; per materiaal worden vervolgens de materiaaleigenschappen ge\"evalueerd alvorens meerdere segmenten worden weggeschreven.

Daarnaast werden alle USE-statements in kaart gebracht in de vorm van een aanroepgraaf. Per USE wordt
genoteerd:
\begin{itemize}
    \item Het type aangeroepen blok: Common Action Block (AB) of External Action Block (EAB).
    \item Of het aangeroepen blok binnen of buiten de migratiescope valt; EAB's vormen vaak natuurlijke
        modulegrenzen omdat ze typisch reeds in PL/I beschikbaar zijn.
\end{itemize}
Het volledige USE-overzicht en bijhorend grafisch schema zijn opgenomen in het analysedossier (zie
Bijlage~\ref{ch:bijlage-views}).

\section{Datamatrix}
\label{sec:re-datamatrix}

Op basis van de READ-, READ EACH-, CREATE-, UPDATE- en DELETE-statements werd een datamatrix opgesteld die
per DB2-tabel weergeeft welke acties plaatsvinden. \autoref{tab:tppbamt-datamatrix} toont het verkorte
overzicht; de gedetailleerde versie inclusief WHERE-clausules en sorteervoorwaarden zit in
Bijlage~\ref{ch:bijlage-views}.

\begin{table}[h]
\centering
\caption[Datamatrix \texttt{AB\_APP\_GBP\_MAT\_BAF}]{Verkorte datamatrix van \texttt{AB\_APP\_GBP\_MAT\_BAF}: gebruik van DB2-tabellen.}
\label{tab:tppbamt-datamatrix}
\begin{tabular}{|l|c|c|c|c|c|p{0.30\linewidth}|}
\hline
\textbf{Tabel} & \textbf{R} & \textbf{RE} & \textbf{C} & \textbf{U} & \textbf{D} & \textbf{Opmerking} \\
\hline
RO & \checkmark & & & & & Roleigenschappen. \\
\hline
\hline
\end{tabular}
\end{table}

\section{Bevindingen}
\label{sec:re-bevindingen}

Uit de analyse komen drie observaties die de aanpak van Hoofdstuk~\ref{ch:poc} sturen:
\begin{enumerate}
    \item De module steunt sterk op Group Views met een vooraf gekende maximumkardinaliteit. Het is
        essentieel dat de PL/I-implementatie dezelfde grenzen respecteert om buffer overflows te vermijden
        (zie \autoref{subsec:groupview-bp}).
    \item Het aantal externe USE-statements is beperkt; het merendeel verwijst naar reeds bestaande PL/I
        EAB's. De impact op de aanroepende generieke modules is daardoor minimaal.
    \item De gehanteerde exitstate-codes wijken licht af van de standaardconventie binnen andere
        modules; in de PL/I-implementatie wordt voor consistentie de standaardconventie.
\end{enumerate}

%Deze observaties bevestigen tegelijk dat de in \autoref{sec:component-analyse} voorgestelde stappen
%voldoende zijn om een module van deze complexiteit te ontsluiten zonder dat een volledige formele
%specificatie nodig is.
